\begin{frame}
\frametitle{Relevanz: Warum 3D-Rekonstruktion in VR?}

\begin{itemize}
  \item Moderne VR-Systeme ermöglichen natürliche Interaktion
  \begin{itemize}
    \item Präzise Hand-, Körper- und Blickerfassung
    \item Stereoskopische Bilddarstellung mit hoher Bildwiederholrate
  \end{itemize}
  
  \item Viele VR-Anwendungen benötigen Bezug zur physischen Umgebung
  \begin{itemize}
    \item Kollaboratives Arbeiten
    \item Simulationsbasiertes Lernen
    \item Räumliche Analysewerkzeuge
  \end{itemize}
  
  \item Virtuelle Räume wirken glaubwürdiger, wenn sie realen Szenen entsprechen
  \begin{itemize}
    \item Erhöhte Immersion und Präsenz
    \item Spontane Erfassung ohne manuelle Modellierung
  \end{itemize}
\end{itemize}

\end{frame}

\begin{frame}
\frametitle{Das Problem: Unzureichende bisherige Lösungen}

\begin{itemize}
  \item Künstlich erstellte Szenen
  \begin{itemize}
    \item Passen schlecht zu realen Umgebungen
    \item Hoher Aufwand für manuelle Modellierung
  \end{itemize}
  
  \item Klassische Photogrammetrie
  \begin{itemize}
    \item Hochauflösende Ergebnisse
    \item Aber: Extensive Bildaufnahmen und lange Verarbeitungszeiten
    \item Ungeeignet für dynamische/interaktive VR-Anwendungen
  \end{itemize}
  
  \item Neue monokulare Rekonstruktionsverfahren existieren
  \begin{itemize}
    \item Nutzen fortlaufende RGB-Videodaten
    \item Ermöglichen schrittweise Geometrierekonstruktion
    \item Problem: Praktische Integration ist schwierig und nicht etabliert
  \end{itemize}
\end{itemize}

\end{frame}

\begin{frame}
\frametitle{Bestehendes Angebot und Forschungslücke}

\begin{itemize}
  \item Monokulare Verfahren passen ideal zu moderner VR-Hardware
  \begin{itemize}
    \item Benötigen nur RGB-Kameras (keine Tiefensensoren)
    \item VR-Headsets haben typischerweise Frontkameras
  \end{itemize}
  
  \item Aber: Verfahren entwickeln sich in isolierten Forschungskontexten
  \begin{itemize}
    \item Eigene Laufzeitumgebungen und Frameworks
    \item Unterschiedliche Datenpipelines und Ausgabeformate
    \item Nicht ausgelegt für interaktive Echtzeit-Systeme
  \end{itemize}
  
  \item Fehlende praktische Integration
  \begin{itemize}
    \item Kein stabiler, strukturierter Datenfluss für VR
    \item Keine standardisierten Schnittstellen zwischen Verfahren
    \item \textbf{Forschungslücke:} Modulare Architektur zur Integration heterogener Verfahren fehlt
  \end{itemize}
\end{itemize}

\end{frame}

\begin{frame}
\frametitle{Zielsetzung dieser Arbeit}

\textbf{Ziel:} Modulares System zur Echtzeit-3D-Rekonstruktion in VR-Umgebung

\bigskip

\begin{itemize}
  \item \textbf{Nicht:} Einzelne Rekonstruktionsmodelle analysieren oder verbessern
  
  \item \textbf{Sondern:} Skalierbare Architektur entwickeln, die
  \begin{itemize}
    \item verschiedene Rekonstruktionsverfahren kapselt
    \item über standardisierte Schnittstelle kommuniziert
    \item mit Unity-basiertem VR-Frontend arbeitet
  \end{itemize}
  
  \item \textbf{Zielumgebung:} Va.Si.Li-Lab der Goethe-Universität Frankfurt
  \begin{itemize}
    \item Bestehende mehrbenutzerfähige VR-Plattform
    \item Modulare, erweiterbare Infrastruktur
    \item Ideale Testumgebung für dynamische 3D-Rekonstruktion
  \end{itemize}
\end{itemize}

\end{frame}

\begin{frame}
\frametitle{Forschungsfrage}

\Large
\centering

\textit{Wie gut eignet sich eine modulare, containerisierte Systemarchitektur zur Integration verschiedener Echtzeit-3D-Rekonstruktionsverfahren in eine bestehende Virtual-Reality-Umgebung?}

\vspace{2cm}

\normalsize
\textbf{Kernfragen der Evaluation:}
\begin{itemize}
  \item Funktioniert die modulare Integration heterogener Verfahren?
  \item Welche Performance-Charakteristiken ergeben sich?
  \item Welche Rekonstruktionsqualität wird erreicht?
  \item Ist das System für praktische VR-Anwendungen geeignet?
\end{itemize}

\end{frame}